\section{Methodology}
\label{sec:methodology}

This section presents the Bi-level Adversarial Reinforcement Learning (Bi-ARL) framework and the BiAT (Bilevel Adversarial Training) framework extension to supervised detectors. We begin with a formal threat model, then describe the Stackelberg game formulation, the bilevel optimization algorithm, and the supervised bilevel transfer.

\subsection{Threat Model}
\label{sec:threat_model}

We consider a \textbf{white-box adaptive adversary} as the primary threat during training, and evaluate under both white-box and gradient-based attacks at test time.

\textbf{Attacker capabilities.}
The attacker has full knowledge of the defender model architecture and weights, and can compute gradients with respect to input features. The attacker's goal is to maximize the false-negative rate (FNR) of the detector by crafting input perturbations $\delta$ with bounded $\ell_\infty$ norm: $\|\delta\|_\infty \leq \epsilon$. Categorical features (e.g., protocol type) are excluded from perturbation.

\textbf{Attacker goal.}
Formally, the attacker solves $\max_{\|\delta\|_\infty \leq \epsilon} \Pr[\hat{y}(x + \delta) = 0 \mid y = 1]$, i.e., maximizing the probability that an attack sample is classified as benign.

\textbf{Defender goal.}
The defender seeks a detection policy $\pi_\mathcal{D}$ that minimizes a composite loss combining false negatives (missed attacks) and false positives (false alarms), even when evaluated against the attacker's best-response strategy.

\textbf{Training vs.\ deployment distinction.}
During \emph{training}, the attacker is white-box and iteratively adapted. At \emph{deployment} (test time), we evaluate against FGSM \cite{goodfellow2015fgsm} and PGD \cite{madry2018pgd} attacks at multiple perturbation budgets to measure the practical robustness of the trained model.

\subsection{Problem Formulation}
\label{sec:formulation}

Let $\mathcal{S}$ denote the state space of tabular traffic-flow features, with dimensionality depending on the dataset (41 for NSL-KDD \cite{tavallaee2009nslkdd}, 42 for UNSW-NB15 \cite{moustafa2015unswnb15}, 78 for CIC-IDS2017 \cite{sharafaldin2018cicids2017} after preprocessing). Two agents interact: Attacker ($\mathcal{A}$) and Defender ($\mathcal{D}$).

At time $t$, the Attacker selects perturbation action $a^{\mathcal{A}}_t \in \mathcal{A}^{\mathcal{A}}$ and the Defender selects classification decision $a^{\mathcal{D}}_t \in \mathcal{A}^{\mathcal{D}}$. We model the interaction as a \textbf{Stackelberg game} \cite{vonStackelberg1952}:

\begin{equation}
    \max_{\theta_{\mathcal{D}}} \; J(\theta_{\mathcal{D}}, \theta_{\mathcal{A}}^*) = \mathbb{E}_{\tau \sim \pi_{\mathcal{D}}, \pi_{\mathcal{A}}^*} \left[ \sum_{t=0}^{T} \gamma^t R_{\mathcal{D}}(s_t, a^{\mathcal{D}}_t, a^{\mathcal{A}*}_t) \right]
    \label{eq:outer}
\end{equation}
\begin{equation}
    \text{s.t.} \quad \theta_{\mathcal{A}}^* = \arg\max_{\theta_{\mathcal{A}}} \; \mathbb{E}_{\tau \sim \pi_{\mathcal{A}}} \left[ \sum_{t=0}^{T} \gamma^t R_{\mathcal{A}}(s_t, a^{\mathcal{D}}_t, a^{\mathcal{A}}_t) \right]
    \label{eq:inner}
\end{equation}

where $\gamma = 0.99$ is the discount factor and $\tau$ denotes a trajectory. Eq.~\eqref{eq:outer} is the outer (defender) problem; Eq.~\eqref{eq:inner} is the inner (attacker best-response) problem. This nesting is what separates our approach from simultaneous-update MARL \cite{lowe2017maddpg, busoniu2008marl}.

\subsection{Reward Functions}

\subsubsection{Attacker Reward}
The Attacker aims to maximize evasion success:
\begin{equation}
    R_{\mathcal{A}} = \begin{cases}
        +2.0 & \text{if FN (attack evades detection)} \\
        -1.0 & \text{if TP (attack detected)}
    \end{cases}
    \label{eq:attacker_reward}
\end{equation}

\subsubsection{Defender Reward}
The Defender minimizes both false negatives and false alarms:
\begin{equation}
    R_{\mathcal{D}} = \begin{cases}
        +1.0 & \text{if TP} \\
        +0.5 & \text{if TN} \\
        -1.0 & \text{if FP} \\
        -2.0 & \text{if FN}
    \end{cases}
    \label{eq:defender_reward}
\end{equation}

The asymmetric penalty ($|r_{\text{FN}}| > |r_{\text{TP}}|$) reflects the higher operational cost of missing an attack relative to generating a false alarm. For UNSW-NB15 and CIC-IDS2017, the false-positive penalty is increased to $-2.5$ to stabilize training under higher class imbalance.

\subsection{PPO-Based Policy Optimization}

Both agents use Proximal Policy Optimization (PPO) \cite{schulman2017ppo}. For policy $\pi_\theta(a|s)$, PPO maximizes:

\begin{equation}
    L^{\text{CLIP}}(\theta) = \mathbb{E}_t \left[ \min \left( r_t(\theta) \hat{A}_t,\; \text{clip}(r_t(\theta), 1{-}\epsilon_{\text{clip}}, 1{+}\epsilon_{\text{clip}}) \hat{A}_t \right) \right]
    \label{eq:ppo}
\end{equation}

where $r_t(\theta) = \pi_\theta(a_t|s_t) / \pi_{\theta_\text{old}}(a_t|s_t)$ is the probability ratio, $\hat{A}_t$ is the advantage estimate from Generalized Advantage Estimation (GAE) \cite{schulman2015gae}:

\begin{equation}
    \hat{A}_t = \sum_{l=0}^{\infty} (\gamma\lambda)^l \delta_{t+l}, \quad \delta_t = r_t + \gamma V(s_{t+1}) - V(s_t)
    \label{eq:gae}
\end{equation}

with $\lambda = 0.95$ and clipping parameter $\epsilon_{\text{clip}} = 0.2$.

\subsection{Bi-level Optimization Algorithm}

\subsubsection{Inner Loop: Attacker Best-Response}
With the defender policy $\pi_\mathcal{D}^{\text{fixed}}$ held constant, the attacker is updated for up to $K_{\text{inner}}$ steps. Convergence is detected by KL divergence:

\begin{equation}
    D_{\text{KL}}(\pi_{\mathcal{A}}^{\text{old}} \| \pi_{\mathcal{A}}^{\text{new}}) < \tau_{\text{KL}}
    \label{eq:kl}
\end{equation}

with $\tau_{\text{KL}} = 0.01$. We use $K_{\text{inner}} = 5$ on NSL-KDD and $K_{\text{inner}} = 1$ on modern datasets (a stabilization choice discussed in Section~\ref{sec:experiments}).

\subsubsection{Outer Loop: Defender Robustification}
Once the attacker reaches $\theta_{\mathcal{A}}^*$, the defender performs one PPO update against this worst-case opponent (frozen in \texttt{eval} mode).

\begin{algorithm}[h]
    \caption{Bi-level Adversarial RL (Bi-ARL)}
    \label{alg:biarl}
    \begin{algorithmic}[1]
        \Require Episodes $N$, inner steps $K_{\text{inner}}$, KL threshold $\tau_{\text{KL}}$
        \Ensure Defender policy $\pi_\mathcal{D}(\cdot|\theta_\mathcal{D})$
        \State Initialize $\theta_\mathcal{D}, \theta_\mathcal{A}$ randomly
        \For{episode $e = 1$ to $N$}
            \State \textbf{Inner Loop: Attacker Best-Response}
            \State $\theta_{\mathcal{A}}^{\text{old}} \leftarrow \theta_{\mathcal{A}}$
            \For{$k = 1$ to $K_{\text{inner}}$}
                \State Sample $\tau \sim \pi_{\mathcal{A}}(\cdot|\theta_{\mathcal{A}}),\; \pi_{\mathcal{D}}^{\text{fixed}}$
                \State $\theta_{\mathcal{A}} \leftarrow \theta_{\mathcal{A}} + \alpha \nabla_{\theta_{\mathcal{A}}} L^{\text{CLIP}}(\theta_{\mathcal{A}})$
                \State Compute $D_{\text{KL}}(\pi_{\mathcal{A}}^{\text{old}} \| \pi_{\mathcal{A}})$
                \If{$D_{\text{KL}} < \tau_{\text{KL}}$}
                    \State \textbf{break} \Comment{Attacker converged}
                \EndIf
            \EndFor
            \State $\theta_{\mathcal{A}}^* \leftarrow \theta_{\mathcal{A}}$ \Comment{Freeze optimal Attacker}
            \State \textbf{Outer Loop: Defender Robustification}
            \State Sample $\tau \sim \pi_{\mathcal{A}}^*(\cdot|\theta_{\mathcal{A}}^*),\; \pi_{\mathcal{D}}(\cdot|\theta_{\mathcal{D}})$
            \State $\theta_{\mathcal{D}} \leftarrow \theta_{\mathcal{D}} + \beta \nabla_{\theta_{\mathcal{D}}} L^{\text{CLIP}}(\theta_{\mathcal{D}})$
        \EndFor
    \end{algorithmic}
\end{algorithm}

\subsection{BiAT Framework: Bilevel Adversarial Training for Supervised Tabular Detectors}
\label{sec:biat}

Bi-ARL establishes the Stackelberg training principle, but the PPO policy network is a relatively weak tabular classifier. The \textbf{BiAT (Bilevel Adversarial Training)} framework transfers the same bilevel idea to supervised detectors with stronger inductive biases for tabular data.

\subsubsection{Motivation}
The ablation results in Section~\ref{sec:experiments} confirm that the bilevel training objective---not just adversarial data augmentation---is responsible for FPR reduction in Bi-ARL. BiAT preserves this objective while replacing the weak RL backbone with MLP and FT-Transformer \cite{gorishniy2021fttransformer} architectures.

\subsubsection{BiAT-MLP}
A three-layer MLP ($256 \to 128 \to 2$, with BatchNorm and Dropout) serves as the baseline supervised detector. The adversarial inner loop uses PGD-style perturbations with $\epsilon = 0.04$.

\subsubsection{BiAT-FTTransformer}
The FT-Transformer \cite{gorishniy2021fttransformer} tokenizes each numerical feature into a $d$-dimensional embedding via a learned linear projection, prepends a \texttt{[CLS]} token, and processes the full sequence through standard transformer encoder layers. Classification uses the \texttt{[CLS]} representation. This architecture's per-feature attention makes it well-suited for tabular IDS data where feature interactions are non-trivial. The adversarial inner loop uses $\epsilon = 0.02$, $\alpha = 0.005$, and 2 PGD steps (calibrated for UNSW-NB15 stability).

\subsubsection{Training Objective}
Both BiAT framework models minimize:
\begin{equation}
    L = (1 - \lambda) L_{\text{clean}} + \lambda L_{\text{adv}}
    \label{eq:biat_loss}
\end{equation}

where $L_{\text{clean}}$ is the cross-entropy on clean samples, $L_{\text{adv}}$ is the cross-entropy on PGD-perturbed samples, and $\lambda \in [0,1]$ controls the adversarial weight. This formulation mirrors the TRADES \cite{zhang2019trades} trade-off between accuracy and robustness. Hyperparameter sensitivity to $\lambda$ is analyzed in Section~\ref{sec:hyperparam}.

\subsubsection{Ablation Chain}
The BiAT framework enables a clean ablation chain:
\begin{center}
    Vanilla FT-Transformer (clean only) $\to$ BiAT-MLP $\to$ BiAT-FTTransformer
\end{center}
This chain isolates the contributions of (a) the stronger backbone and (b) the bilevel adversarial training objective.

\subsection{Theoretical Analysis}
\label{sec:theory}

This section formalizes the approximation quality of the finite-step inner loop and motivates the stabilization mechanisms introduced in our framework.

\subsubsection{Finite-Step Inner Loop Approximation}

Let $\pi_{\mathcal{A}}^*(\theta_{\mathcal{D}})$ denote the true best-response attacker policy against a fixed defender $\theta_{\mathcal{D}}$, and let $\pi_{\mathcal{A}}^{(K)}(\theta_{\mathcal{D}})$ denote the attacker policy obtained after $K$ inner-loop PPO update steps. We characterize the approximation error below.

\begin{proposition}[Inner-Loop Approximation Error Bound]
\label{prop:inner_approx}
Assume the attacker's PPO objective $J_{\mathcal{A}}(\theta_{\mathcal{A}} | \theta_{\mathcal{D}})$ is $L$-smooth and $\mu$-strongly concave in $\theta_{\mathcal{A}}$ (within the trust region enforced by clipping). After $K$ inner-loop steps with learning rate $\alpha \leq 1/L$, the policy gap satisfies:
\begin{equation}
    \left\| \pi_{\mathcal{A}}^{(K)} - \pi_{\mathcal{A}}^* \right\|_{\mathrm{TV}} \leq \left(1 - \frac{\mu}{L}\right)^{K/2} \cdot C_0
    \label{eq:inner_bound}
\end{equation}
where $C_0 = \left\| \pi_{\mathcal{A}}^{(0)} - \pi_{\mathcal{A}}^* \right\|_{\mathrm{TV}}$ is the initial gap and $\|\cdot\|_{\mathrm{TV}}$ denotes total variation distance.
\end{proposition}

\textit{Proof sketch.} Under $\mu$-strong concavity (guaranteed locally by PPO clipping) and $L$-smoothness, gradient ascent converges linearly. Converting the parameter-space convergence rate $\|\theta^{(K)} - \theta^*\| \leq (1-\mu/L)^K \|\theta^{(0)} - \theta^*\|$ to policy space via the mean-value theorem for neural network parameterizations yields the stated bound \cite{sinha2018bilevel}. $\square$

\textbf{Practical implication.} With $K=5$ and $\mu/L \approx 0.15$ (estimated from training curves), the bound gives $(0.85)^{2.5} \approx 0.68$, indicating that about 32\% of the initial policy gap is closed per outer-loop iteration. This aligns with our empirical finding that $K_{\text{inner}}=5$ provides the best accuracy--robustness trade-off (see Section~\ref{sec:hyperparam}).

\begin{proposition}[Defender Sub-optimality Bound]
\label{prop:defender_subopt}
Let $V_{\mathcal{D}}^*$ be the defender's value under the true Stackelberg equilibrium, and let $V_{\mathcal{D}}^{(K)}$ be the defender's value when trained against the $K$-step approximate attacker. If the defender's value function is $L_V$-Lipschitz with respect to the attacker policy, then:
\begin{equation}
    V_{\mathcal{D}}^* - V_{\mathcal{D}}^{(K)} \leq L_V \cdot \left(1 - \frac{\mu}{L}\right)^{K/2} \cdot C_0
    \label{eq:defender_bound}
\end{equation}
\end{proposition}

\textit{Proof sketch.} This follows directly from Proposition~\ref{prop:inner_approx} and the Lipschitz assumption: the defender's performance degradation is bounded by the attacker's deviation from best-response, scaled by the sensitivity constant $L_V$. $\square$

\textbf{Remark.} The exponential decay in $K$ justifies using a moderate $K_{\text{inner}}$ (e.g., 5) rather than running the inner loop to convergence, as the marginal benefit diminishes rapidly while the computational cost grows linearly.

\subsubsection{Inner-Loop Stabilization via Entropy-Regularized Stackelberg Games}
\label{sec:stabilization}

In practice, the strong concavity assumption may not hold globally, leading to inner-loop instability (as observed in our UNSW-NB15 experiments). We address this by formulating the inner loop as an \textbf{entropy-regularized} Stackelberg game:

\begin{equation}
    \max_{\pi_{\mathcal{A}}} \; J_{\mathcal{A}}(\pi_{\mathcal{A}} | \theta_{\mathcal{D}}) + \beta(t) \cdot \mathcal{H}(\pi_{\mathcal{A}})
    \label{eq:entropy_reg}
\end{equation}

where $\mathcal{H}(\pi_{\mathcal{A}}) = -\mathbb{E}[\log \pi_{\mathcal{A}}]$ is the policy entropy and $\beta(t)$ is a time-dependent coefficient with exponential decay schedule $\beta(t) = \max(\beta_{\min},\; \beta_0 \cdot \gamma_\beta^t)$.

The entropy term serves a dual purpose:
\begin{enumerate}
    \item \textbf{Regularization}: The augmented objective $J + \beta \mathcal{H}$ is $\beta$-strongly concave in the policy simplex \cite{geist2019entropy}, ensuring the inner loop satisfies the conditions of Proposition~\ref{prop:inner_approx} even when the original $J_{\mathcal{A}}$ is not strongly concave.
    \item \textbf{Exploration}: Preventing premature attacker policy collapse improves the quality of the best-response approximation, as the attacker explores a broader range of adversarial strategies.
\end{enumerate}

Additionally, we employ an \textbf{attacker warm-up schedule}: during the first $\rho$ fraction of training episodes ($\rho = 0.2$), the inner loop uses $K_{\text{inner}} = 1$, linearly increasing to the full value thereafter. This prevents early-stage training instability caused by the attacker over-optimizing against an untrained defender.

